\chapter*{Einleitung zur dritten Lerneinheit}
\addcontentsline{toc}{chapter}{Einleitung zur dritten Lerneinheit}

Die vorliegende Lerneinheit beschäftigt sich ausschließlich mit Funktionen, wobei zu Beginn vor allem auf die unterschiedlichen Arten der Darstellung einer Funktion eingegangen wird.

Es werden dann nur sogenannte reelle Funktionen mit einer veränderlichen Größe (die meistens $x$ heißt) betrachtet, die in der Regel durch eine Formel beschreibbar sind (zumindest stückweise). Zur grafischen Darstellung bedient man sich des üblichen $x$-$y$-Koordinatensystems.

Zunächst werden die wichtigsten Eigenschaften von Funktionen untersucht. Insbesondere spielen die Schnittpunkte mit den Koordinatenachsen und das »Verhalten im Unendlichen« eine Rolle.

Der nächste Abschnitt ist speziellen Funktionen gewidmet. Es werden ganz-rationale, gebrochen rationale, irrationale und transzendente Funktionen behandelt. In jeder dieser Klassen werden die für die Anwendungen wichtigsten Funktionen betrachtet, d.\,h. mit ihren Eigenschaften hinsichtlich Definitions- und Wertebereich, Stetigkeit etc. vorgestellt. Auf die Behandlung der vor allem in technischen Zusammenhängen bedeutenden trigonometrischen Funktionen wird verzichtet.

Im letzten Abschnitt geht es um die Differenzialrechnung. Dabei wird der Ableitungsbegriff ohne formale Grenzwertbetrachtung vorgestellt. Auch auf die Technik des Differenzierens wird so weit eingegangen, dass die »gängigsten« Funktionen abgeleitet werden können. Den Abschluss bilden die Themen Kurvendiskussion sowie Extremwertaufgaben als typisches Anwendungsfeld der Differenzialrechnung.

Die verwendeten Beispiele werden -- wenn möglich -- dem Bereich der ökonomischen Funktionen entnommen, da auch in dieser Lerneinheit die betriebswirtschaftliche Perspektive im Vordergrund stehen soll.

Die Lerneinheit beinhaltet einige Übungsaufgaben, die Lösungen kann man sich auf den letzten Seiten ansehen. Weitere Übungsaufgaben sind auch im Internet (z.\,B. auf »youtube«) in zahlreichen Mathematik-Angeboten zu finden.
